\documentclass[12pt,a4paper]{article}
\usepackage[utf8]{inputenc}
\usepackage[T1]{fontenc}
\usepackage{geometry}
\usepackage{titlesec}
\usepackage{hyperref}
\usepackage{graphicx}
\usepackage{fancyhdr}
\usepackage{enumitem}
\usepackage{tabularx}
\usepackage{float}
\geometry{left=2.5cm,right=2.5cm,top=2.5cm,bottom=2.5cm}

\title{\textbf{ChurnZero: Modern Customer Success Methodologies}}
\author{Summary of Webinar Insights}
\date{\today}

\begin{document}

\maketitle

\section{The One-Slide Value Review}
\textit{A Modern Alternative to Traditional QBRs}

Traditional Quarterly Business Reviews (QBRs) often suffer from being too long (30+ slides), backward-looking, and lacking strategic momentum. The "One-Slide Value Review" condenses the entire strategic conversation into a single, high-impact slide divided into four quadrants.

\subsection{The Four Quadrants}
\begin{enumerate}
    \item \textbf{Expected Value:} Re-anchors the conversation on the customer's original business goals.
    \begin{itemize}
        \item \textit{Key Question:} Why did they buy? What is the primary outcome they are seeking?
    \end{itemize}
    
    \item \textbf{Value Delivered:} Focuses on impact rather than activity.
    \begin{itemize}
        \item \textit{Shift:} Instead of "We closed 5 tickets," say "We saved you 10 hours of downtime."
    \end{itemize}
    
    \item \textbf{Proof of Value:} meaningful data visualizing progress.
    \begin{itemize}
        \item \textit{Visuals:} Show the Baseline $\rightarrow$ Current State $\rightarrow$ Target Goal.
    \end{itemize}
    
    \item \textbf{Next Steps (Risks \& Asks):}
    \begin{itemize}
        \item \textbf{Relentless Focus:} What needs to happen in the next 90 days?
        \item \textbf{Executive Asks:} Explicitly request help to unblock risks.
    \end{itemize}
\end{enumerate}

\subsection{Strategic Best Practices}
\begin{itemize}
    \item \textbf{Pre-Reads:} Send data-heavy charts and usage reports via email \textit{before} the meeting. Don't waste meeting time reviewing them.
    \item \textbf{Executive Presence:} Keep the meeting to 30 minutes to respect executive time.
    \item \textbf{Alignment:} Ensure the "Champion" and "Executive Sponsor" are both present.
\end{itemize}

\section{Engagement AI}
ChurnZero's Engagement AI helps Customer Success teams identify risk and understand relationship health through automated analysis of communications.

\subsection{Core Features}
\begin{itemize}
    \item \textbf{Sentiment Analysis:} Scans emails, survey responses, and meeting notes to automatically tag interactions as Positive, Neutral, or Negative.
    \item \textbf{Topic Tracking:} Identifies recurring themes (e.g., "Bugs," "Renewal," "Pricing") to spot trends without manual tagging.
    \item \textbf{Relationship Score:} A real-time health signal derived from the aggregate sentiment of recent interactions.
    \item \textbf{Aggregated Reporting:} A dashboard view that visualizes communication trends over time, helping identifying drifting accounts.
\end{itemize}

\section{CRO Expectations from Customer Success}
Insights from Chief Revenue Officers on the evolving role of CS in an efficiency-first economy.

\subsection{Key Expectations}
\begin{itemize}
    \item \textbf{Commercial Mindset:} CS is not just "Support." It is a revenue engine responsible for Net Revenue Retention (NRR).
    \item \textbf{Efficient Growth:} Doing more with less. Leveraging AI to handle simple inquiries so humans can focus on high-value strategy.
    \item \textbf{Action-Based Health Scores:} Moving away from "Subjective Health" (e.g., "I feel they are happy") to "Objective Health" based on actions (e.g., "They adopted Feature X").
    \item \textbf{Alignment:} Deep integration between Sales and CS to ensure the "Promise" made by Sales is the "Reality" delivered by CS.
\end{itemize}

\end{document}
